% GENERATED FILE. Edit canonical lessons, then run npm run generate:books.
% canonical-source-hash: fnv1a64:38532c0a33742205
% canonical-lessons: PA-C127-kinara, PA-C127-jharna, PA-C127-gufa, PA-C127-vadi, PA-C127-tibba

\chapter{Shore, Waterfall, Cave, Valley, Sand hill}
\label{ch:pa-shore-waterfall-cave-valley-sand-hill}
\hlchaptermodality{127}

\begin{chapteropening}
\textbf{By the end of this chapter:} \emph{I can say a shore, a waterfall, a cave, a valley and a sand hill.}

Complete the last lesson of chapter 127: I can say a shore, a waterfall, a cave, a valley and a sand hill.
\end{chapteropening}

\section[\texorpdfstring{\pa{ਕਿਨਾਰਾ} (kinārā)}{ਕਿਨਾਰਾ (kinārā)}]{\pa{ਕਿਨਾਰਾ} (kinārā) — a shore, a bank}

\label{lesson:PA-C127-kinara}



\begin{quote}
Before the new one: say the Punjabi for a bud, then the Punjabi for an island.
\end{quote}

\subsection*{You'll want to know: \pa{ਕਿਨਾਰਾ}}
\textbf{\pa{ਕਿਨਾਰਾ}} — \emph{kinārā} — \textquotedblleft{}a shore, a bank\textquotedblright{}.

\begin{cousinweb}
\textbf{\pa{ਕਿਨਾਰਾ}} (\emph{kinārā}) is a loan, usually traced to Persian \emph{kinār}, \textquotedblleft{}edge, shore\textquotedblright{}.
\end{cousinweb}

\begin{grammarlens}[title={What you've built}]
One more word for this chapter.
\end{grammarlens}

\subsection*{Your turn}
\begin{itemize}
\raggedright
  \item \emph{Say it:} \emph{kinārā}
  \item \emph{Say it:} \emph{kinārā}, once more
  \item \emph{Say it:} say \emph{kinārā}
  \item \emph{Recall:} say the Punjabi for a bud, then the Punjabi for an island, then say \emph{kinārā} again
\end{itemize}

\begin{tcolorbox}[breakable,colback=teal!4,colframe=teal!35!black,title={Before you move on}]
What does \pa{ਕਿਨਾਰਾ} mean? (\textquotedblleft{}A shore, a bank\textquotedblright{}.) Say it once more.
\end{tcolorbox}
\section[\texorpdfstring{\pa{ਝਰਨਾ} (jharnā)}{ਝਰਨਾ (jharnā)}]{\pa{ਝਰਨਾ} (jharnā) — a waterfall}

\label{lesson:PA-C127-jharna}



\begin{quote}
Before the new one: say the Punjabi for an island, then the Punjabi for a shore.
\end{quote}

\subsection*{You'll want to know: \pa{ਝਰਨਾ}}
\textbf{\pa{ਝਰਨਾ}} — \emph{jharnā} — \textquotedblleft{}a waterfall\textquotedblright{}.

\begin{cousinweb}
\textbf{\pa{ਝਰਨਾ}} (\emph{jharnā}) belongs to the root \emph{jhar}, \textquotedblleft{}flow down\textquotedblright{}, but CDIAL does not list it in the sense \textquotedblleft{}waterfall\textquotedblright{}. It gives Hindi \emph{jharī}, \textquotedblleft{}waterfall\textquotedblright{}, under Sanskrit \emph{jhara}, \textquotedblleft{}waterfall\textquotedblright{} (CDIAL 5343), and a Punjabi \emph{jharnā} only as \textquotedblleft{}grating in a waterway\textquotedblright{}, under \emph{jharaṇa}, as in \emph{jharaṇodaka}, \textquotedblleft{}water from a cascade\textquotedblright{} (CDIAL 5344).
\end{cousinweb}

\begin{grammarlens}[title={What you've built}]
One more word for this chapter.
\end{grammarlens}

\subsection*{Your turn}
\begin{itemize}
\raggedright
  \item \emph{Say it:} \emph{jharnā}
  \item \emph{Say it:} \emph{jharnā}, once more
  \item \emph{Say it:} say \emph{jharnā}
  \item \emph{Recall:} say the Punjabi for an island, then the Punjabi for a shore, then say \emph{jharnā} again
\end{itemize}

\begin{tcolorbox}[breakable,colback=teal!4,colframe=teal!35!black,title={Before you move on}]
What does \pa{ਝਰਨਾ} mean? (\textquotedblleft{}A waterfall\textquotedblright{}.) Say it once more.
\end{tcolorbox}
\section[\texorpdfstring{\pa{ਗੁਫ਼ਾ} (gufā)}{ਗੁਫ਼ਾ (gufā)}]{\pa{ਗੁਫ਼ਾ} (gufā) — a cave}

\label{lesson:PA-C127-gufa}



\begin{quote}
Before the new one: say the Punjabi for a shore, then the Punjabi for a waterfall.
\end{quote}

\subsection*{You'll want to know: \pa{ਗੁਫ਼ਾ}}
\textbf{\pa{ਗੁਫ਼ਾ}} — \emph{gufā} — \textquotedblleft{}a cave\textquotedblright{}.

\begin{cousinweb}
CDIAL lists \textbf{\pa{ਗੁਫ਼ਾ}} (\emph{gufā}), in the spelling \emph{guphā}, \textquotedblleft{}cave\textquotedblright{}, under a reconstructed \emph{*gupphā}, \textquotedblleft{}cave\textquotedblright{} (CDIAL 4204). CDIAL calls its own further link with Sanskrit \emph{guhā}, \textquotedblleft{}cave\textquotedblright{}, very doubtful.
\end{cousinweb}

\begin{grammarlens}[title={What you've built}]
One more word for this chapter.
\end{grammarlens}

\subsection*{Your turn}
\begin{itemize}
\raggedright
  \item \emph{Say it:} \emph{gufā}
  \item \emph{Say it:} \emph{gufā}, once more
  \item \emph{Say it:} say \emph{gufā}
  \item \emph{Recall:} say the Punjabi for a shore, then the Punjabi for a waterfall, then say \emph{gufā} again
\end{itemize}

\begin{tcolorbox}[breakable,colback=teal!4,colframe=teal!35!black,title={Before you move on}]
What does \pa{ਗੁਫ਼ਾ} mean? (\textquotedblleft{}A cave\textquotedblright{}.) Say it once more.
\end{tcolorbox}
\section[\texorpdfstring{\pa{ਵਾਦੀ} (vādī)}{ਵਾਦੀ (vādī)}]{\pa{ਵਾਦੀ} (vādī) — a valley}

\label{lesson:PA-C127-vadi}



\begin{quote}
Before the new one: say the Punjabi for a waterfall, then the Punjabi for a cave.
\end{quote}

\subsection*{You'll want to know: \pa{ਵਾਦੀ}}
\textbf{\pa{ਵਾਦੀ}} — \emph{vādī} — \textquotedblleft{}a valley\textquotedblright{}.

\begin{cousinweb}
\textbf{\pa{ਵਾਦੀ}} (\emph{vādī}) is a loan from Arabic \emph{wādī}, \textquotedblleft{}valley, river bed\textquotedblright{}, which reached Punjabi through Persian. Its Arabic root is \emph{w-d-y} (Lane).
\end{cousinweb}

\begin{grammarlens}[title={What you've built}]
One more word for this chapter.
\end{grammarlens}

\subsection*{Your turn}
\begin{itemize}
\raggedright
  \item \emph{Say it:} \emph{vādī}
  \item \emph{Say it:} \emph{vādī}, once more
  \item \emph{Say it:} say \emph{vādī}
  \item \emph{Recall:} say the Punjabi for a waterfall, then the Punjabi for a cave, then say \emph{vādī} again
\end{itemize}

\begin{tcolorbox}[breakable,colback=teal!4,colframe=teal!35!black,title={Before you move on}]
What does \pa{ਵਾਦੀ} mean? (\textquotedblleft{}A valley\textquotedblright{}.) Say it once more.
\end{tcolorbox}
\section[\texorpdfstring{\pa{ਟਿੱਬਾ} (ṭibbā)}{ਟਿੱਬਾ (ṭibbā)}]{\pa{ਟਿੱਬਾ} (ṭibbā) — a sand hill}

\label{lesson:PA-C127-tibba}



\begin{quote}
Before the new one: say the Punjabi for a cave, then the Punjabi for a valley.
\end{quote}

\subsection*{You'll want to know: \pa{ਟਿੱਬਾ}}
\textbf{\pa{ਟਿੱਬਾ}} — \emph{ṭibbā} — \textquotedblleft{}a sand hill\textquotedblright{}.

\begin{cousinweb}
\textbf{\pa{ਟਿੱਬਾ}} (\emph{ṭibbā}) is inherited. CDIAL lists it, \textquotedblleft{}mound, sandbank, hill\textquotedblright{}, under \emph{*ṭibba}, one of the variant forms it groups under \emph{*ṭappa}, \textquotedblleft{}hill\textquotedblright{} (CDIAL 5446).
\end{cousinweb}

\begin{grammarlens}[title={What you've built}]
One more word for this chapter.
\end{grammarlens}

\subsection*{Your turn}
\begin{itemize}
\raggedright
  \item \emph{Say it:} \emph{ṭibbā}
  \item \emph{Say it:} \emph{ṭibbā}, once more
  \item \emph{Say it:} say \emph{ṭibbā}
  \item \emph{Recall:} say the Punjabi for a cave, then the Punjabi for a valley, then say \emph{ṭibbā} again
\end{itemize}

\begin{tcolorbox}[breakable,colback=teal!4,colframe=teal!35!black,title={Before you move on}]
What does \pa{ਟਿੱਬਾ} mean? (\textquotedblleft{}A sand hill\textquotedblright{}.) Say it once more.
\end{tcolorbox}
